\documentclass[11pt]{article}
\usepackage{amsmath,amssymb,amsthm,mathtools}
\usepackage{geometry}
\geometry{margin=1in}
\usepackage{enumitem}
\usepackage{booktabs}

\newtheorem{theorem}{Theorem}
\newtheorem{lemma}{Lemma}
\newtheorem{proposition}{Proposition}
\newtheorem{definition}{Definition}
\newtheorem{warning}{Warning}

\title{Step 148: Source-Compatible Weight Normalization Audit}
\author{RATCHET / Membrane RH Construction Notes}
\date{}

\begin{document}
\maketitle

\section{Purpose}
Step 146 showed that the unweighted residual zero ledger is not trace-class-certified on the Burnol/Sonine residual carrier. Step 147 supplied a positive weighted ledger
\[
  K_R^\omega=(D_R^\omega)^*D_R^\omega,
  \qquad
  (D_R^\omega h)_{\rho,k}=\sqrt{\omega_{\rho,k}}\,q(\rho)\,\langle h,\Pi_RY^a_{\rho,k}\rangle,
\]
with
\[
  q(\rho)=\Re\rho-\frac12.
\]
The present step audits the remaining compatibility condition:
\[
  C_{\rm src}^\omega
  :=\sup_N \operatorname{tr}(F_N^\Omega K_R^\omega)<\infty,
\]
or more generally
\[
  \operatorname{tr}(F_N^\Omega K_R^\omega)=o(\Lambda_N^\Omega).
\]
The verdict is that this is not a free normalization choice. It is a load-bearing source-audit theorem.

\section{Objects}
Let $H_R$ be the completed Burnol/Sonine residual carrier from Step 145. Let $G_R\succeq0$ denote the completed residual metric. Let
\[
  F_N^\Omega\succeq \Lambda_N^\Omega G_R - E_N
\]
be the lifted $\Omega$-compatible restricted source frame. In the idealized no-defect statement we write
\[
  F_N^\Omega\succeq \Lambda_N^\Omega G_R,
  \qquad \Lambda_N^\Omega\to\infty.
\]
Let $K_R^\omega\ge0$ be the fixed positive trace-class weighted residual ledger.

\section{Weighted source-trace formula}
For the rank-one/evaluator representation
\[
  K_R^\omega
  =\sum_{(\rho,k)} \omega_{\rho,k}q(\rho)^2
     |\Pi_RY^a_{\rho,k}\rangle\langle \Pi_RY^a_{\rho,k}|,
\]
whenever the sum is trace-class and $F$ is positive bounded on the relevant form domain,
\[
\boxed{
  \operatorname{tr}(F K_R^\omega)
  =\sum_{(\rho,k)}\omega_{\rho,k}q(\rho)^2
    \langle F\Pi_RY^a_{\rho,k},\Pi_RY^a_{\rho,k}\rangle.
}
\]
Thus source compatibility is a source-evaluator response-envelope statement, not merely a height-weight summability statement.

\section{No-free-normalization theorem}
\begin{theorem}[No free source-compatible normalization]
Assume $F_N\succeq\Lambda_NG_R$, $\Lambda_N\to\infty$, and $K\ge0$ is trace-class. Then
\[
  \operatorname{tr}(F_NK)
  \ge
  \Lambda_N\operatorname{tr}(G_RK).
\]
Consequently, if
\[
  \operatorname{tr}(F_NK)=o(\Lambda_N),
\]
then
\[
  \operatorname{tr}(G_RK)=0.
\]
In particular, a uniformly bounded source budget
\[
  \sup_N\operatorname{tr}(F_NK)<\infty
\]
forces completed weighted residual collapse.
\end{theorem}

\begin{proof}
Since $F_N-\Lambda_NG_R\succeq0$ and $K\ge0$ is trace-class,
\[
  \operatorname{tr}((F_N-\Lambda_NG_R)K)\ge0.
\]
Hence
\[
  \operatorname{tr}(F_NK)\ge\Lambda_N\operatorname{tr}(G_RK).
\]
Dividing by $\Lambda_N$ gives the conclusion.
\end{proof}

\section{Interpretation}
This theorem cuts both ways. It is exactly the completed squeeze mechanism, but it also says that source compatibility cannot be purchased by declaring faster-decaying fixed positive weights if the residual ledger has nonzero mass in the source metric. For a fixed separating ledger, bounded source exposure is already a collapse-level certificate.

\section{Envelope-normalized schedules}
Let dyadic height shells be denoted by $\mathcal Z_j$, with
\[
  n_j=\#\mathcal Z_j,
  \qquad
  E_j=\sup_{z\in\mathcal Z_j}\|\Pi_RY_z\|^2.
\]
If one also has a shell source exposure envelope
\[
  U_j=\sup_N\sup_{z\in\mathcal Z_j}
  \frac{\langle F_N\Pi_RY_z,\Pi_RY_z\rangle}{\|\Pi_RY_z\|^2},
\]
then a formal envelope schedule
\[
  \omega_z=\frac{2^{-j}}{(1+n_j)(1+E_j)(1+U_j)},
  \qquad z\in\mathcal Z_j,
\]
controls the shellwise source trace. However, when $F_N\succeq\Lambda_NG_R$ on a fixed nonzero direction and $\Lambda_N\to\infty$, the envelope $U_j$ is infinite unless that direction is absent from the completed residual ledger. Therefore this construction is an exposure-control lemma, not a proof of completed residual collapse.

\section{Moving-weight warning}
A window-dependent weight schedule such as
\[
  \omega_z^{(N)}\sim (\Lambda_N^\Omega)^{-1}\omega_z
\]
can artificially keep
\[
  \operatorname{tr}(F_NK_R^{\omega^{(N)}})
\]
bounded. But then the ledger is not fixed. The result has status
\[
  \texttt{moving\_window\_support\_only}.
\]
It cannot replace the completed fixed/exhaustive ledger needed for the residual theorem.

\section{Source-compatible squeeze}
With tail terms restored, the accepted completed squeeze has the form
\[
  \operatorname{tr}(G_RK_R^\omega)
  \le
  \frac{\operatorname{tr}(F_N^\Omega K_R^\omega)}{\Lambda_N^\Omega}
  +\operatorname{tr}(T_NK_R^\omega)
  +\text{declared defects}.
\]
Thus the completed residual ledger collapses if
\[
  \frac{\operatorname{tr}(F_N^\Omega K_R^\omega)}{\Lambda_N^\Omega}\to0,
\]
\[
  \operatorname{tr}(T_NK_R^\omega)\to0,
\]
and all declared defects vanish.

\section{Verdict}
The Step 147 weighted ledger is trace-class and separating under its envelope assumptions. Step 148 shows that source compatibility is the next genuine audit theorem:
\[
\boxed{
  \operatorname{tr}(F_N^\Omega K_R^\omega)=o(\Lambda_N^\Omega).
}
\]
This cannot be obtained merely by renormalizing weights without either proving collapse, forfeiting fixed-ledger status, or adding a new source-budget law.

\section{Next step}
Step 149 should be the source audit budget theorem or obstruction:
\[
\boxed{
  \operatorname{tr}(F_N^\Omega K_R^\omega)=o(\Lambda_N^\Omega)
  \quad ?
}
\]
If the theorem fails, the route must identify the source-exposed residual sector as a new $\Xi$-residual.

\end{document}
